\chapter{Abelian Gauge Theory and QED}

\section{Local phase symmetry}

Begin with a Dirac field and the global transformation
\[
  \psi\to e^{-\ii q\alpha}\psi.
\]
The free Dirac Lagrangian
\[
  \Lag=\bar\psi(\ii\gamma^\mu\p_\mu-m)\psi
\]
is invariant if \(\alpha\) is constant.  If \(\alpha=\alpha(x)\), then
\[
  \p_\mu(e^{-\ii q\alpha(x)}\psi)
  =
  e^{-\ii q\alpha}
  \left(\p_\mu\psi-\ii q(\p_\mu\alpha)\psi\right).
\]
The extra term spoils invariance.

Introduce a gauge field \(A_\mu\) and define
\[
  D_\mu=\p_\mu+\ii q A_\mu.
\]
Demand
\[
  D_\mu\psi\to e^{-\ii q\alpha}D_\mu\psi.
\]
This holds if
\[
  A_\mu\to A_\mu+\p_\mu\alpha.
\]
Indeed,
\[
  D_\mu'\psi'
  =
  (\p_\mu+\ii qA_\mu+\ii q\p_\mu\alpha)
  (e^{-\ii q\alpha}\psi)
  =
  e^{-\ii q\alpha}(\p_\mu+\ii qA_\mu)\psi.
\]

\section{Field strength}

The electromagnetic field strength is
\[
  F_{\mu\nu}=\p_\mu A_\nu-\p_\nu A_\mu.
\]
It is gauge invariant:
\[
  F_{\mu\nu}\to
  \p_\mu(A_\nu+\p_\nu\alpha)
  -
  \p_\nu(A_\mu+\p_\mu\alpha)
  =
  F_{\mu\nu},
\]
because mixed partial derivatives commute.

The gauge-field Lagrangian is
\[
  \Lag_{\text{gauge}}=-\frac14F_{\mu\nu}F^{\mu\nu}.
\]
Together with matter,
\[
  \Lag_{\text{QED}}
  =
  \bar\psi(\ii\gamma^\mu D_\mu-m)\psi
  -\frac14F_{\mu\nu}F^{\mu\nu}.
\]
Expanding the covariant derivative,
\[
  \bar\psi\ii\gamma^\mu(\p_\mu+\ii qA_\mu)\psi
  =
  \bar\psi\ii\gamma^\mu\p_\mu\psi
  -q\bar\psi\gamma^\mu A_\mu\psi.
\]
The interaction current is
\[
  j^\mu=q\bar\psi\gamma^\mu\psi.
\]

\section{Maxwell equations from the action}

Vary \(A_\mu\) in
\[
  \Lag=-\frac14F_{\rho\sigma}F^{\rho\sigma}-j^\mu A_\mu.
\]
Since
\[
  \delta F_{\rho\sigma}=\p_\rho\delta A_\sigma-\p_\sigma\delta A_\rho,
\]
and \(F^{\rho\sigma}\) is antisymmetric,
\[
  \delta\left(-\frac14F_{\rho\sigma}F^{\rho\sigma}\right)
  =
  -\frac12F^{\rho\sigma}\delta F_{\rho\sigma}
  =
  -F^{\rho\sigma}\p_\rho\delta A_\sigma.
\]
Integrating by parts,
\[
  \delta S
  =
  \int\dd^4x
  \left[
  \p_\rho F^{\rho\sigma}-j^\sigma
  \right]\delta A_\sigma.
\]
Therefore
\[
  \p_\rho F^{\rho\sigma}=j^\sigma.
\]
The remaining homogeneous Maxwell equations follow from the definition of \(F_{\mu\nu}\):
\[
  \p_\lambda F_{\mu\nu}
  +\p_\mu F_{\nu\lambda}
  +\p_\nu F_{\lambda\mu}=0.
\]

\section{Gauge fixing and the photon propagator}

The kinetic operator for \(A_\mu\) is not invertible until gauge redundancy is handled.  Add the covariant gauge-fixing term
\[
  \Lag_{\text{gf}}=-\frac1{2\xi}(\p_\mu A^\mu)^2.
\]
In Feynman gauge \(\xi=1\), the quadratic gauge-field Lagrangian becomes
\[
  \Lag_2=\frac12 A_\mu \eta^{\mu\nu}\Boxop A_\nu
\]
after integration by parts.  The photon propagator is
\[
  \frac{-\ii\eta_{\mu\nu}}{p^2+\ii\epsilon}.
\]
In Abelian QED, Faddeev--Popov ghosts decouple because the gauge transformation of \(\p_\mu A^\mu\) is independent of \(A_\mu\).  In non-Abelian theory they do not decouple.

\section{QED Feynman rules}

The main QED rules are:
\[
  \text{fermion propagator:}\quad
  \frac{\ii(\slashed p+m)}{p^2-m^2+\ii\epsilon},
\]
\[
  \text{photon propagator in Feynman gauge:}\quad
  \frac{-\ii\eta_{\mu\nu}}{p^2+\ii\epsilon},
\]
\[
  \text{vertex:}\quad -\ii q\gamma^\mu.
\]
Arrows on fermion lines keep track of charge flow.  Closed fermion loops acquire an extra minus sign because anticommuting fields must be moved past one another.

\section{Ward identity}

Gauge symmetry implies relations among amplitudes.  A basic form is that replacing an external photon polarization \(\epsilon_\mu(k)\) by its momentum \(k_\mu\) gives zero:
\[
  k_\mu\mathcal M^\mu=0.
\]
For a fermion line emitting a photon,
\[
  k_\mu\gamma^\mu=\slashed k.
\]
If the adjacent fermion momenta are \(p\) and \(p+k\), then
\[
  \slashed k=(\slashed p+\slashed k-m)-(\slashed p-m).
\]
When placed between on-shell spinors satisfying
\[
  (\slashed p-m)u(p)=0,
  \qquad
  \bar u(p+k)(\slashed p+\slashed k-m)=0,
\]
the expression vanishes.  This elementary cancellation is the seed of the Ward identities.

\section*{Exercises}
\addcontentsline{toc}{section}{Exercises}

\begin{exercise}
Show explicitly that \(F_{\mu\nu}\) is unchanged by \(A_\mu\to A_\mu+\p_\mu\alpha\).
\begin{answer}
The extra terms are \(\p_\mu\p_\nu\alpha-\p_\nu\p_\mu\alpha=0\).
\end{answer}
\end{exercise}

\begin{exercise}
Derive the QED vertex factor from the interaction term \(-q\bar\psi\gamma^\mu A_\mu\psi\).
\begin{answer}
The interaction enters the path integral as \(\exp[\ii\int\Lag_I]\), so one use of the interaction term gives \(-\ii q\gamma^\mu\).
\end{answer}
\end{exercise}
